\newglossaryentry{N}{name={$N$}, description={the set of nodes in a graph}}
\newglossaryentry{E}{name={$E$}, description={the set of edges in a graph}}
\newglossaryentry{G}{name={$G$}, description={a graph}}
\newglossaryentry{AllNodes}{name={$\AllNodes$}, description={the alphabet of graph nodes}}
\newglossaryentry{AllLabels}{name={$\AllLabels$}, description={the alphabet of graph edge labels}}
\newglossaryentry{rrt}{name={$\rrt$}, description={re-rooting operation}}
\newglossaryentry{laa}{name={$\la{a}$}, description={an edge with label $a$}}
\newglossaryentry{las}{name={$\las$}, description={a path, possibly empty, leading to}}
\newglossaryentry{isoeqf}{name={$\isoeq[f]$}, description={isomorphic to, under bijective function $f$}}
\newglossaryentry{coversf}{name={$\covers[f]$}, description={covering, under graph isomorphism $f$}}
\newglossaryentry{gmatf}{name={$\gmat[f]$}, description={matching, under graph isomorphism $f$}}
\newglossaryentry{gdaf}{name={$\gda[f]$}, description={de-aliased from, under total function $f$}}
\newglossaryentry{coversgda}{name={$\covers\gda$}, description={covering de-aliased graph of}}
\newglossaryentry{isosubgda}{name={$\isosub\gda$}, description={covered by de-aliased graph of}}
\newglossaryentry{isa}{name={$\isa$}, description={subclass of, is-a}}
\newglossaryentry{subty}{name={$\subty$}, description={subtype of, preorder}}
\newglossaryentry{bgconj}{name={$\bgconj$}, description={graph conjunction}}
\newglossaryentry{bgdisj}{name={$\bgdisj$}, description={graph disjunction}}
\newglossaryentry{AllClasses}{name={$\AllClasses$}, description={the alphabet of classes}}
\newglossaryentry{AllValues}{name={$\AllValues$}, description={the alphabet of values}}
\newglossaryentry{DFS}{name={\DFS}, description={depth-first search}}

\chapter{Pattern Graphs and Match Relation}\label{ch:graph}
\minitoc

We first give the formal definitions of graphs, operations, morphisms and relations.
We begin with the basic layout graphs consisting of only nodes and edges, and the match relation w.r.t. to a given preorder.
Then, we add type assignments to the layout graphs to form the pattern graphs, and describe how to encode program states as graphs.
We further expand the theory to match conjunction and disjunction, discuss the restriction and provide our resolution.

\section{Layout Graphs}

Let $\AllNodes$ be the set of all nodes, which can be any unique entities, usually associated with types, values and objects, and
$\AllLabels$ the set of all labels, which are usually variable names. We define the layout graphs on these elements.

\begin{Definition}[Layout Graph]
A layout graph $\mgls{G} = \vect{\mgls{N}, \mgls{E}, p}$ is an edge-labeled, rooted and directed graph, where
\begin{enumerate}
\item $\Nodes(G) = N \subset \mgls{AllNodes}$ is the set of nodes,

\item $\Edges(G) = E: N \to \mgls{AllLabels} \to N$ is a partial function, in infix notation, representing the edges, mapping the source node and the label to the target node, and

\item $\Root(G) = p \in N$ is the root, and all the nodes are connected to $p$, $\forall u \in N \bullet \exists \overline{a} \bullet p \overline{E a} = u$.\footnote{%
The overline $\overline{a}$ notation represents a series $a_1,\dots,a_n$ of items.
However, an variable $E$ bound in the context, if overlined, repeats itself rather than expands into a series,
thus $\overline{Ea}$ expands to $Ea_1Ea_2 \dots Ea_n$.}
\end{enumerate}
\end{Definition}

We formulate the edges in a layout graph as a \emph{curried} function to imply that the outgoing edges from a node must have unique labels.
We write, more intuitively, $u \mathrel{\mgls{laa}} v \in E$ to denote an edge $u E a = v$.
Also, $u \mathrel{\mgls{las}} v \in E$ denotes there exists a path from $u$ to $v$, including single-node paths.

A layout graph represents the layout of an object, where the root becomes the origin.
Other nodes of a layout graph are the origins of object components,
and each component is represented by the subgraph connecting to its origin.

\begin{Definition}[Re-rooting]\label{def:reroot}
Let $G = \vect{N,E,p}$ be a layout graph, and $p' \in N$.
We say $G' = \vect{N',E',p'}$ is the subgraph re-rooted at node $p'$,
denoted as $G' = G \mathbin{\mgls{rrt}} p'$, where
\begin{enumerate}
\item the nodes are connected to $p'$,
\[
N' = \sett{v \mid \exists p' \las v \in E},
\]
\item edges are restricted to the connected nodes, $E' = E \proj{N'}$.
\end{enumerate}
\end{Definition}

Because there may be circular paths going back to the root,
re-rooting cannot guarantee a \emph{proper} subgraph.
Therefore, it does not necessarily imply a proper reduction.
We must explicitly remove all the incoming edges to the root, and examine the residue subgraph.
We define the subtraction to remove some edges from a layout graph.

\begin{Definition}[Subtraction]\label{def:gesub}
Let $G = \vect{N,E,p}$ be a layout graph, and $X \subseteq E$ the edges to exclude.
$G' = \vect{N',E',p}$ is the subtraction graph of $X$ from $G$,
denoted as $G' = G \gsub X$, where
\begin{enumerate}
\item the nodes are connected via $I = E \gsub X$ to the root,
\[
N' = \sett{v \mid \exists p \las v \in I},
\]
\item edges are restricted to the connected nodes, $E' = I\proj{N'}$.
\end{enumerate}
\end{Definition}

Since morphisms are functions on nodes, the subtraction gives us a way to break cycles in a graph
for a proper reduction on the number of nodes through re-rooting.

\section{Match Relation}

The structural equivalence of two layout graphs is formalized as an isomorphism,
which is a bijection preserving all the outgoing edges of a node.

\begin{Definition}[Graph Isomorphism]\label{def:giso}
Let $G = \vect{N,E,p}$ and $G' = \vect{N',E',p'}$ be two layout graphs.
$G$ is isomorphic to $G'$, denoted as $G \mathrel{\mgls{isoeqf}} G'$,
if there exists a bijective function $f: N \to N'$, such that
\begin{enumerate}
\item root maps to root, $f(p) = p'$, and

\item edges are preserved, $u \la{a} v \in E \iff f(u) \la{a} f(v) \in E'$.
\end{enumerate}
If $G$ is isomorphic to a root-preserving subgraph $G''$ of $G'$, i.e., $\Root(G'') = \Root(G')$,
we say $G'$ \emph{covers} $G$, denoted as  $G' \mathrel{\mgls{coversf}} G$ or $G \isosub[f] G'$.
\end{Definition}

\begin{figure}[t]
	\centering
	\begin{tikzpicture}[default,->,thick,x={(45:45pt)}, y={(160:40pt)}, z={(115:12pt)}]
		\node[root] (r) at (0,0,0) {};
		\node[vertex] (a) at (0,-2,0) {};
		\node[vertex] (b) at (1,-1,0) {};
		\node[vertex] (c) at (2,0,0) {};
		\node[vertex] (d) at (2,-2,0) {};

		\draw (r) .. controls +(-75:1) and +(105:1) .. node[below] {$x$} (a);
		\draw (r) .. controls +(-30:0.75) and +(150:0.75) .. node[below,pos=0.7] {$y$} (b);
		\draw (r) .. controls +(15:0.75) and +(150:0.5) .. node[above,pos=0.7] {$z$}(c);
		\draw (b) .. controls +(-150:1) and +(30:0.75) .. node[right,pos=0.8] {$x$} (a);
		\draw (b) .. controls +(-30:0.75) and +(150:0.75) .. node[below,pos=0.7] {$y$} (d);
		\draw (c) .. controls +(-30:1) and +(30:1) .. node[right] {$z$} (c);

		\node[root,Maroon] (r1) at (0,0,1) {};
		\node[vertex,Maroon] (a1) at (0,-2,1) {};
		\node[vertex,Maroon] (b1) at (1,-1,1) {};

		\draw[Maroon] (r1) .. controls +(-75:1) and +(105:1) .. (a1);
		\draw[Maroon] (r1) .. controls +(-30:0.75) and +(150:0.75) .. (b1);
		\draw[Maroon] (b1) .. controls +(-150:1) and +(30:0.75) .. (a1);

		\node[root,Orchid] (r2) at (0,0,2) {};
		\node[vertex,Orchid] (b2) at (1,-1,2) {};
		\node[vertex,Orchid] (c2) at (2,0,2) {};
		\node[vertex,Orchid] (d2) at (2,-2,2) {};

		\draw[Orchid] (r2) .. controls +(-30:0.75) and +(150:0.75) .. (b2);
		\draw[Orchid] (r2) .. controls +(15:0.75) and +(150:0.5) .. (c2);
		\draw[Orchid] (b2) .. controls +(-30:0.75) and +(150:0.75) .. (d2);
		\node at ($(c2)+(0.5,0,0)$) {};
		\useasboundingbox;
		\draw[Orchid] (c2) .. controls +(-30:1) and +(30:1) .. (c2);

		\draw[thin,-,MidnightBlue] (0,0,-2) -- (r);
		\draw[thin,-,MidnightBlue] (r.center) -- (r1);
		\draw[thin,-,MidnightBlue] (r1.center) -- (r2);
		\draw[thin,-,MidnightBlue] (r2.center) -- (0,0,4);

		\draw[thin,-,Periwinkle] (a.center) -- (a1);
		\draw[thin,-,Periwinkle] (b.center) -- (b1);
		\draw[thin,-,Periwinkle] (b1.center) -- (b2);
		\draw[thin,-,Periwinkle] (c.center) -- (c2);
		\draw[thin,-,Periwinkle] (d.center) -- (d2);
	\end{tikzpicture}
	\caption{Subgraph isomorphism and covering between layout graphs}\label{fig:subiso}
\end{figure}

\autoref{fig:subiso} illustrates the covering between layout graphs.
An object not only has the layout, but also has a type.
We associate the type of an object to the root of its layout graph through a type-assignment $\tau$.
Suppose there is a preorder $\subty$ defined on these types, e.g., the subclass or subset relation.
The type-assignment is there because nodes of graphs are unique, it's inconvenient to reason about ordering between nodes directly.
We define the match relation $\gmat$ between two layout graphs.
\begin{Definition}[Graph Match]\label{def:gmat}
Let $G$ and $G'$ be two layout graphs, and $\mgls{subty}$ a preorder on $\ran(\tau)$.
$G'$ matches $G$ w.r.t. $\subty$ under isomorphism $f$,
denoted as $G' \gmatord[f]{\subty} G$, or simply $G' \mathrel{\mgls{gmatf}} G$ without ambiguity,
if $G' \covers[f] G$, such that
the preorder exists between the types of any pair of corresponding nodes, i.e., $\forall u \in N \bullet \tau(f(u)) \subty \tau(u)$.
\end{Definition}

The graph match is a cover relation restricted to the given preorder between the types of nodes.
It is trivial to verify that the graph match $\gmat$ relation is reflexive.
We can also prove that $\gmat$ is transitive.

\begin{Theorem}
The graph match $\gmat$ relation is a preorder.
\end{Theorem}
\begin{proof}
Suppose $G' = \vect{N',E',p'} \gmat[f] G = \vect{N,E,p}$
and $G'' = \vect{N'',E'',p''} \gmat[f'] G'$.
We can verify
\begin{enumerate}
\item $f'(f(p)) = p''$,

\item $u \la{a} v \in E \implies f'(f(u)) \la{a} f'(f(v)) \in E''$, and

\item $\forall u \in N \bullet \tau(f'(f(u))) \subty \tau(u)$.
\end{enumerate}
Hence $G'' \gmat[f' \circ f] G$. The $\gmat$ relation is transitive.
\end{proof}

Obviously, if a layout graph $G$ is matched by another graph $G'$, any re-rooted subgraph of $G$ must also be matched
by the subgraph of $G'$ re-rooted at the corresponding node.

\begin{Theorem}\label{thm:submatch}
If $G' \gmat[f] G = \vect{N,E,p}$, $q \in N$ and $X \subseteq E$,
then $G' \rrt f(q) \gmat (G \gsub X) \rrt q$.
\end{Theorem}
\begin{proof}
By restricting $f$ to the nodes of $(G \gsub X) \rrt q$.
\end{proof}

\begin{Corollary}\label{thm:match-unique}
If there are two morphisms $f_1$ and $f_2$ such that $G' \gmat G$ under $f_1$ and $f_2$,
then $f_1$ = $f_2$, i.e., the match morphism is unique.
\end{Corollary}
\begin{proof}
By induction on the number of nodes. We remove all the incoming edges to the root, and apply \autoref{thm:submatch}.
\end{proof}

On the other hand, by first removing all the incoming edges to the root,
we can then re-root the residue subgraph to the adjacent nodes of the root,
to reduce the match problem to the subgraphs.
By combining the morphisms for these re-rooted subgraphs, if they exist, and adding back the root to the combined morphism,
we are able to obtain the morphism for the original layout graph, if it exists.

\begin{Theorem}\label{thm:match-reduction}
Given two layout graphs $G = \vect{N,E,p}$ and $G' = \vect{N',E',p'}$,
Let $G^\ast = G \gsub \sett{u \la{a} v \in E \mid v = p}$
be the residue subgraph of $G$ with all the incoming edges to the root removed.
If
\begin{enumerate}
\item for all the outgoing edges from the root $p\la{a}q \in E$,
\[
\exists q' \suchthat p'\la{a}q' \in E' \wedge (p = q \vee G' \rrt q' \gmat[f_q] G^\ast \rrt q),
\]
\item and for all the incoming edges to the root $u\la{b}p \in E$,
\[
\exists q \suchthat f_q(u)\la{b}p' \in E',
\]
\end{enumerate}
then $G' \gmat[f] G$, where $f = \left(\bigcup_q f_q\right) \cup \sett{p \mapsto p'}$.
\end{Theorem}
\begin{proof}
Obviously, by the given conditions, the addition of $\sett{p \mapsto p'}$ preserves the incoming and outgoing edges of the root.
By the definition of layout graphs, all the nodes are connected to the root, therefore any node in a layout graph must be either the root,
or in one of the subgraphs re-rooted at the nodes adjacent to the root.
Hence all we need to prove is that $\bigcup_q f_q$ exists, i.e., if $u \in \dom(f_q) \cap \dom(f_r)$ then $f_q(u) = f_r(u)$.
This is true by \autoref{thm:submatch} and \autoref{thm:match-unique}, that any (sub)match morphism is unique.
\end{proof}

Although finding the subgraph isomorphism is NP-complete in general, it can be efficiently done for rooted graphs by a simple depth-first search (\gls{DFS}),
led to by \autoref{thm:match-reduction}.
We can determine whether there is a match between two layout graphs by trying to construct the match morphism during the {\DFS}.
If the construction fails, we can conclude that the match morphism does not exist.

\begin{algorithm}[t]
\caption{Construction of Match Morphism}\label{algo:match-dfs-cons}
\DontPrintSemicolon
\SetKwFunction{DfsMatch}{dfs\_match}
\KwIn{preorder:~$\subty$, type-assignment:~$\tau$, layout~graphs:~$G$ and $G'$.}
\KwOut{match morphism:~$f$ if $G' \gmat[f] G$.}
\KwEx{$\Mismatch$.}
\Begin{
  $\vect{N,E,p} = G, \vect{N',E',p'} = G'$\;
  $f \gets \emptyset$\;
  \DfsMatch{$E, p, E', p'$}\;
}
\medskip
\Fn{\DfsMatch{$E, p, E', p'$}}{
  \eIf{$p \in \dom(f)$}{
    \lIRaise{$\Mismatch$}{%
      $f(p) \neq p'$
    }
  }{
    \lIRaise{$\Mismatch$}{%
      $\tau(p') \nsubty \tau(p)$
    }
    \lIRaise{$\Mismatch$}{%
      $p' \in \ran(f)$
    }
    $f \gets f \cup \sett{p \mapsto p'}$\;
    \ForAll{$p \la{a} q \in E$}{
      \lIRaise{$\Mismatch$}{%
        $p' \la{a} q' \notin E'$
      }
      \DfsMatch{$E, q, E', q'$}\;
    }
  }
}
\end{algorithm}

\autoref{algo:match-dfs-cons} gives the precise construction procedure.
It takes the preorder and type-assignment on the nodes and two layout graphs as the input, and
output a morphism mapping between the nodes of the two graphs, if the match can be established, otherwise an exception is raised.
The procedure initializes the morphism to an empty map, and starts traversing both layout graphs synchronously from their corresponding roots.
The {\DfsMatch} function performs the {\DFS}. It checks the \mbox{1-1} mapping and the preorder relation between the two roots, and
accumulates the mapping to the morphism.
Then the function checks, w.r.t. the edge labels, the correspondence of the adjacent nodes, and recurs on each pair of them.

\section{Pattern Graphs}

A pattern graph $P = \vect{G, T}$ is a layout graph $P$ with a map $T$ of type-assignment,
which associates each node with a type. A type can be
a usual class,
a subclass,
a set of values and
particularly a singleton set of one value.
The preorder $\subty$ on the node types of such pattern graphs is exactly the subtype relation between two types.
In principle, a type $t'$ is a subtype of another type $t$ if all objects or values of $t'$ are elements of $t$,
i.e., $\Elem(t') \subseteq \Elem(t)$.

\begin{table}[t]
\centering
\caption{Subtyping Rules}\label{tab:subtyping}
\begin{mathpar}
\inferrulel[(Reflexive)]
  {c \in \AllClasses}
  {c \subty c}
\and
\inferrulel[(Subclass)]
  {c,d \in \AllClasses \\ \Env \vdash d \mathrel{\mgls{isa}} c}
  {\Env \vdash d \subty c}
\and
\inferrulel[(Subset)]
  {W \subseteq V \subseteq \AllValues}
  {W \subty V}
\and
\inferrulel[(Value-Type)]
  {V \subseteq \AllValues \\ t \in \AllClasses \\ \Env \vdash \forall v \in V \suchthat v:t}
  {\Env \vdash V \subty t},
\end{mathpar}
\end{table}

\begin{Definition}[Subtype]\label{def:subtype}
A type is either a class $c \in \mgls{AllClasses}$ or a set of values $V \subset \mgls{AllValues}$.
The subtype relation $\subty$ is inferred from the typing environment $\Env$ by the subtyping rules given in \autoref{tab:subtyping},
where $\isa$ is the subclass relation.
\end{Definition}

\begin{Theorem}
The subtype relation $\subty$  in \autoref{def:subtype} is a preorder.
\end{Theorem}
\begin{proof}
All we need is that the subclass relation $\isa$ being unidirectional and transistive, and so it is.
\end{proof}

\begin{Definition}[Pattern Match]
Given two pattern graphs $P = \vect{G, T}$, $P' = \vect{G', T'}$ and typing environment $\Env$,
$P'$ matches $P$ w.r.t. preorder $\subty$ on node types, denoted as $P' \gmatord{\subty} P$,
if $G' \gmatord{\subty} G$, where $\tau(u') \subty \tau(u) \qdef \Env \vdash T'(u') \subty T(u)$.
\end{Definition}

An object corresponds to a node in a pattern graph, with the root representing the main object,
and the edges leading to component objects.
In most cases, an object must terminate at some value nodes, which are primitives in a programming language, such as integers and booleans.
It is common to use the value itself as the node. This formulation complicates the graph model by having concrete and multiple types of nodes.
To keep the nodes pure and abstract, we associate a node $p$ with its value $\nvalue(p)$ through a run-time function.
We then store this value as a singleton set of one value when constructing the type map of a pattern graph.
The real type of the node can be obtained from the value, since any terminus does not start a structure, and must have a known type in the type system.

The state $\Sta$ of a running program can be represented by a layout graph $H$,
together with the above mentioned $\nvalue$ function and the run-time $\atype$ information,
where $H$ consists of a linked list of scope nodes as a stack,
with outgoing edges as local variables to the roots of the objects~\cite{ke2009graph}.
Therefore, a pattern graph can be extracted from a state by re-rooting $H$ at the root $p$ of the object.
The map of type assignment is constructed from $\atype(p)$,
which returns the actual type of the object.

\begin{Definition}[Pattern Extraction]
A pattern graph $P = \vect{G,T}$ is extracted from a state $\Sta = \vect{H,\nvalue,\atype}$ at node $p$,
denoted as $P = \Sta \rrt p$, where $G = H \rrt p$, and for all $u \in \Nodes(G)$,
\[
T(u) = \begin{cases}
\atype(u) & \mbox{if } u \notin \dom(\nvalue), \\
\sett{\nvalue(u)} & \mbox{if } u \in \dom(\nvalue).
\end{cases}
\]
\end{Definition}

\section{Match Conjunction and Disjunction}

When we use pattern matching, it is natural to think about matching one object against multiple patterns,
and take the conjunction or disjunction of the individual matches.
When taking the conjunction of multiple patterns, we want to know if their overlapping parts are consistent.
This observation can tell us if it is possible to have some object matching all these patterns simultaneously.
On the other hand, when taking the disjunction of multiple patterns, the overlapping parts give us the common component of the individual matches.
This component carries the properties that we can ensure no matter which branch succeeds.

To check whether multiple pattern graphs are consistent, we try to construct a union graph that matches these layouts.
The construction is divided into two parts. First, we make sure that an isomorphism is possible from each layout to the union.
This tells whether the patterns are structurally consistent.
Second, we make sure that there exists at least one type that is a subtype of all the types associated with the corresponding nodes of the patterns,
better there exists an infimum. This tells whether the patterns are consistent in the type system.

Aliasing is an inherit feature of object references, if a node is a common target of two edges, it can be mapped to two nodes of different layouts separately.
However, there is no way for this node to be mapped to the two nodes of the union graph simultaneously.
We need to address the aliasing issue when we use the union graph to deal with match conjunction.

\begin{Definition}[De-aliasing]
Let $G = \vect{N,E,p}$ and $G' = \vect{N',E',p'}$ be two layout graphs.
$G$ is a de-aliased graph of $G'$,
denoted as $G \mathrel{\mgls{gdaf}} G'$,
if there exists a total function $f: N \to N'$, such that
\begin{enumerate}
\item root is preserved, $f(p) = p'$,

\item edges are preserved, $u \la{a} v \in E \implies f(u) \la{a} f(v) \in E'$ and
$u' \la{a} v' \in E' \implies \exists u\la{a}v \in E \suchthat (u',v') = (f(u),f(v))$.
\end{enumerate}

For a layout graph $G''$,
if $\exists G \gda G' \suchthat G'' \covers G$,
we say $G''$ covers some de-aliased graph of $G'$,
denoted as $G'' \mathrel{\mgls{coversgda}} G'$.
If $\exists G \gda G' \suchthat G'' \isosub G$,
we say $G''$ is covered by some de-aliased graph of $G'$,
denoted as $G'' \mathrel{\mgls{isosubgda}} G'$.
\end{Definition}

\begin{Definition}[Union]
$\co{G}$ is a union of multiple layout graphs $\overline{G}$,
denoted as $\co{G} \isoeq[f] \bigcup\overline{G}$,
if $\overline{G \isosub[f] \co{G}}$, and
for all $G'$ such that $\overline{G \isosub G'}$,
we have $\co{G} \isosub\gda G'$.

The union is the minimum layout graph that covers all the given layouts, and cannot be further de-aliased.
\end{Definition}

\begin{Theorem}
Given multiple layout graphs $\overline{G}$,
if there exists a layout graph $G'$ that covers all of them,
there exists the union of $\overline{G}$.
\end{Theorem}
\begin{proof}
Since we have $\overline{G \isosub[f] G'}$, we can construct the union $\co{G}$ from $\overline{f}$ and $\overline{G}$.
The nodes of $\co{G}$ are mapped from $\overline{\dom(f)}$, and combined only when necessary.
If a node $u_i$ in $\dom(f_i)$ has the set of incoming paths intersecting with that of a node $u_j$ in $\dom(f_j)$,
both $u_i$ and $u_j$ should be mapped to one node~\cite{hoare1999trace}.
Otherwise, they should be mapped to different nodes.
The existence of $\overline{f}$ guarantees the set of incoming paths of a node in one layout will not lead to more than one node in another layout.
Such construction splits the target nodes whenever possible, ensuring that the result cannot be further de-aliased.
\end{proof}

\begin{algorithm}[t]
\caption{Construction of Union}\label{algo:union-dfs-cons}
\DontPrintSemicolon
\SetKwFunction{DfsUnion}{dfs\_union}
\KwIn{layout graphs:~$\overline{G}$.}
\KwOut{morphisms to union:~$\overline{f}$.}
\KwEx{$\NoUnion$.}
\Begin{
  $\overline{\vect{N,E,p} = G}$\;
  $\overline{f \gets \emptyset}$\;
  \DfsUnion{$\overline{E},\overline{p}$}\;
}
\medskip
\Fn{\DfsUnion{$\overline{E_i},\overline{p_i}$}}{
  $\overline{k} = \sett{k \in \overline{i} \mid p_k \in \dom(f_k)}$\;
  $\overline{c} = \overline{i} \gsub \overline{k}$\;
  $\overline{u = f_k(p_k)}$\;
  \lIRaise{$\NoUnion$}{$|\sett{\overline{u}}| > 1$}
  \If{$\overline{c} \neq \emptyset$}{
    \eIf{$\overline{u} \neq \emptyset$}{
      \lIRaise{$\NoUnion$}{$u_1 \in \bigcup\overline{\ran(f_c)}$}
      $\co{p} \gets u_1$
    }{
      $\co{p} \gets \text{new node}$
    }
    $\overline{f_c \gets f_c \cup \sett{p_c \mapsto \co{p}}}$\;
    \ForAll{$a,\overline{j} = \sett{j \in \overline{i} \mid p_j\la{a}q_j \in E_j}$}{
      \DfsUnion{$\overline{E_j},\overline{q_j}$}
    }
  }
}
\end{algorithm}

As illustrated in \autoref{algo:union-dfs-cons}, we try to construct the morphisms $\overline{f}$ from each graph to the union
by traversing the graphs simultaneously in {\DFS}, while checking for conflicts.
In each iteration of the {\DFS}, we collect the nodes $\overline{p_i}$ that have a common incoming path only in those graphs $\overline{G_i}$,
where $\overline{G_i}$ is a subset of $\overline{G}$.
These nodes $\overline{p_i}$ should be mapped to the same node in the union.
We divide $\overline{p_i}$ into two sets, one set $\overline{p_k}$ are previously visited from other paths, the other set $\overline{p_c}$ are not yet visited.
There are two conflicts to check for.
First, the visited nodes must all be mapped to one node $u$.
Second, this $u$ must not be an existing target from those graphs $\overline{G_c}$ where the unvisited nodes reside,
which means there are two paths leading to one node in $\overline{G_k}$, however leading to two nodes in $\overline{G_c}$.
If no conflicts are detected, we assign $u$, or a new node if $u$ does not exist, as the mapping target of those unvisited nodes $\overline{p_c}$.

\begin{Definition}[Intersection]
$\ov{G}$ is an intersection of multiple layout graphs $\overline{G}$,
denoted as $\ov{G} \isoeq[f \circ g] \bigcap\overline{G}$,
if $\overline{G \covers[f]\gda[g] \ov{G}}$, and
for all $G'$ such that $\overline{G \covers\gda G'}$,
we have $\ov{G} \covers\gda G'$.

The intersection is the maximum layout graph that all the given layouts cover some de-aliased version of it.
\end{Definition}

\begin{Theorem}
Given multiple layout graphs $\overline{G}$,
there always exists the intersection of $\overline{G}$.
\end{Theorem}
\begin{proof}
We can construct the intersection as follows.
First, the intersection begins with collecting those incoming paths common in all $\overline{G}$.
Then, if any two paths reach the same node in any one of $\overline{G}$, we combine them into one set of incoming paths,
otherwise, we keep them as two separate sets.
Finally, each set of incoming paths is mapped to a node of the intersection.
Similar to the union, such construction splits the target nodes whenever possible, ensuring that the result cannot be further de-aliased.
\end{proof}

\begin{algorithm}[t]
\caption{Construction of Intersection}\label{algo:inter-dfs-cons}
\DontPrintSemicolon
\SetKwFunction{DfsInter}{dfs\_inter}
\KwIn{layout graphs:~$\overline{G}$.}
\KwOut{morphisms to intersection:~$\overline{f}$.}
\Begin{
  $\overline{\vect{N,E,p} = G}$\;
  $\overline{f \gets \emptyset}$\;
  \DfsInter{$\overline{E},\overline{p}$}
}
\medskip
\Fn{\DfsInter{$\overline{E},\overline{p}$}}{
  $\overline{k} = \sett{k \mid p_k \in \dom(f_k)}$\;
  $\overline{u = f_k(p_k)}$\;
  \eIf{$\overline{u} \neq \emptyset$}{
    $\ov{p} \gets u_1$
  }{
    $\ov{p} \gets \text{new node}$
  }
  $\overline{f \gets f \cup \sett{p \mapsto \ov{p}}}$\;
  \If{$|\overline{u}| \neq |\overline{p}|$}{
    \ForAll{$a \suchthat \overline{p\la{a}q \in E}$}{
      \DfsInter{$\overline{E},\overline{q}$}\;
    }
  }
}
\end{algorithm}

As illustrated in \autoref{algo:inter-dfs-cons}, the construction of intersection is simpler.
We again traverse the graphs $\overline{G}$ simultaneously in {\DFS}.
In each iteration, we collect only those node $\overline{p}$ with a common incoming path in all the given layouts.
If some of the nodes $\overline{p_k}$ are previously visited, they may have been mapped to multiple nodes $\overline{u}$.
We take one of the targets $u_1$ to unify the mapping target of all these nodes.
If all $\overline{p}$ are unvisited, we map all of them to a new node.

The union and intersection give us the common structures of multiple layout graphs,
however without respect to the preorder between node types.
If there is an object matching all the given layouts $\overline{G}$,
it must contain a subgraph $G'$ that can be de-aliased to the union of $\overline{G}$.
Each node $u'$ of $G'$ corresponds to one node $\co{u}$ of the union in each component match of the conjunction.
This $\co{u}$ must therefore be of a subtype of all the nodes mapped to it from $\overline{G}$.
We can assume a better type of $\co{u}$ and thus $u'$, if there exists an infimum type of all the nodes mapped to $\co{u}$.

\begin{Definition}[Match Conjunction]
$G'$ matches the conjunction of layout graphs $\overline{G}$, w.r.t. preorder $\subty$,
denoted as $G' \gmatord{\subty} \mgls{bgconj}\overline{G}$,
if $\overline{G' \gmatord{\subty} G}$.
\end{Definition}

\begin{algorithm}[t]
\caption{Construction of Match Conjunction}\label{algo:conj-dfs-cons}
\DontPrintSemicolon
\KwIn{typing environment:~$\Env$, pattern graphs:~$\overline{P}$.}
\KwOut{morphisms to union:~$\overline{f}$ and type~assignment:~$\co{T}$.}
\KwEx{$\NoUnion,\NoConj$.}
\Begin{
  $\overline{\vect{G,T} = P}$\;
  \autoref{algo:union-dfs-cons}: $\overline{G} \to \overline{f}$\;
  $\co{T} \gets \emptyset$\;
  $\co{N} = \bigcup\overline{\ran(f)}$\;
  \ForAll{$\co{u} \in \co{N}$}{
    $\overline{k} = \sett{k \mid \co{u} \in \ran(f_k)}$\;
    $\overline{t_k = T_k(f_k^{-1}(\co{u}))}$\;
    \lIRaise{$\NoConj$}{$\Env \vdash \nexists t \suchthat \overline{t \subty t_k}$}
    $\Env \vdash \co{t} = \inf(\overline{t_k})$\;
    $\co{T} \gets \co{T} \cup \sett{\co{u} \mapsto \co{t}}$
  }
}
\end{algorithm}

\autoref{algo:conj-dfs-cons} constructs the type assignment of the union based on the given patterns.
We first obtain the union morphisms from the layouts by calling \autoref{algo:union-dfs-cons}.
Then, for each node $\co{u}$ in the union, we collect all the layouts having nodes mapped to $\co{u}$.
These layouts correspond to morphisms $\overline{f_k}$ and type assignments $\overline{T_k}$.
Thus, we are able to look up the types $\overline{t_k}$ associated with the nodes mapped to $\co{u}$, and
check for the legitimacy of subtyping $\overline{t_k}$.
If the typing environment does not permit such subtyping, we must conclude that the conjunction of these patterns are not possible.
If the subtyping is possible, we try to find the infimum $\co{t}$ of $\overline{t_k}$ and associate it with $\co{u}$.
The type of $\co{u}$ is $\NoType$ if there exists no infimum.

On the contrary, for match disjunction, an object is to match only one of the multiple patterns.
This is always possible because neither structural nor subtyping conflicts can occur for only one graph.
Any object that matches one of the given layouts $\overline{G}$ must contain a subgraph $G'$
which is a de-aliased graph of the intersection of $\overline{G}$.
Each node $u'$ of $G'$ is mapped to the corresponding node $\ov{u}$ of the intersection
possibly via any component match of the disjunction.
This $\ov{u}$ must therefore be of a supertype of all the nodes mapped to it,
to cope with all the possible individual matches.
Even if such a supertype exists, it is possibly not unique.
We can only associate a type with $\ov{u}$ when there is a supremum type of all the nodes mapped to $\ov{u}$.

\begin{Definition}[Match Disjunction]
$G'$ matches the disjunction of layout graphs $\overline{G}$, w.r.t. preorder $\subty$,
denoted as $G' \gmatord{\subty} \mgls{bgdisj}\overline{G}$,
if there exists $G'' \in \overline{G}$ such that $G' \gmatord{\subty} G''$.
\end{Definition}

\begin{algorithm}[t]
\caption{Construction of Match Disjunction}\label{algo:disj-dfs-cons}
\DontPrintSemicolon
\KwIn{typing environment:~$\Env$, pattern graphs:~$\overline{P}$.}
\KwOut{morphisms to intersection:~$\overline{f}$ and type~assignment:~$\ov{T}$.}
\Begin{
  $\overline{\vect{G,T} = P}$\;
  \autoref{algo:inter-dfs-cons}: $\overline{G} \to \overline{f}$\;
  $\ov{T} \gets \emptyset$\;
  $\ov{N} = \bigcup\overline{\ran(f)}$\;
  \ForAll{$\ov{u} \in \ov{N}$}{
    $\overline{t} = \bigcup\overline{\sett{T(u) \mid f(u) = \ov{u}}}$\;
    $\Env \vdash \ov{t} = \sup(\overline{t})$\;
    $\ov{T} \gets \ov{T} \cup \sett{\ov{u} \mapsto \ov{t}}$
  }
}
\end{algorithm}

\autoref{algo:disj-dfs-cons} constructs the type assignment of the intersection based on the given patterns.
We first obtain the intersection morphisms from the layouts by calling \autoref{algo:inter-dfs-cons}.
Then, for each node $\ov{u}$ in the intersection, we find the nodes mapped to $\ov{u}$ in all the patterns.
Note that a morphism to the intersection is not injective, there can be more than one nodes mapped to $\ov{u}$.
Next, we look up the types $\overline{t}$ associated with these nodes.
Finally, we try to find the supremum $\ov{t}$ of $\overline{t}$ and associate it with $\ov{u}$.
The type of $\ov{u}$ is $\NoType$ if there exists no supremum.

\endinput